\subsection{Adsorption from Solution}
Adsorption from solution involves the substitution of solvent molecules by solute molecules on the surface.
\textbf{Gibbs Adsorption Equation (for solutions):}
\begin{equation}
    \Gamma = -\frac{C}{RT} \frac{d\gamma}{dC}
\end{equation}
\textbf{Surface Excess amount ($n^\sigma$):}
\begin{equation}
    n_B^\sigma = n_B^s - (n_B^s + n_A^s) x_B^b
\end{equation}
where $n^s$ is total amount at the surface, $x^b$ is mole fraction in bulk.

\textbf{Rebindia's Rule (Polarity):} Substances are adsorbed most strongly at interfaces where the difference in polarity between the two phases is greatest. Adsorption occurs in the direction that balances the polarity difference.

\subsection{Adsorption of Electrolytes}
\textbf{Basic Rules:}
\begin{enumerate}
    \item \textbf{Hydration Effect:} For ions of the same valence, those with larger crystal radii (smaller hydrated radii) are adsorbed more strongly (\textbf{Lyotropic Series}):
        \begin{itemize}
            \item $Li^+ < Na^+ < K^+ < Rb^+ < Cs^+$
            \item $Mg^{2+} < Ca^{2+} < Sr^{2+} < Ba^{2+}$
            \item $Cl^- < Br^- < NO_3^- < I^- < CNS^-$
        \end{itemize}
    \item \textbf{Valence Effect:} Higher valence ions are adsorbed more strongly: $K^+ < Ca^{2+} < Al^{3+} < Th^{4+}$.
    \item \textbf{Structural Priority:} Priority is given to ions that can fit into the crystal lattice of the adsorbent.
\end{enumerate}

\subsection{Ion Exchange}
Process of exchange of ions between an electrolyte solution and a solid (ionite).
\textbf{Classification of Ion Exchangers:}
\begin{itemize}
    \item \textbf{Cationic:} Exchange positive ions ($H^+$, metal cations). Functional groups: $-\text{SO}_3H$, $-\text{COOH}$, $-\text{OH}$.
    \item \textbf{Anionic:} Exchange negative ions ($OH^-$, anions). Functional groups: $-\text{NH}_2$, $=\text{NH}$, $\equiv\text{N}$.
    \item \textbf{Amphoteric:} Can exchange both.
\end{itemize}
\textbf{Structure:}
\begin{itemize}
    \item \textbf{Inorganic:} Natural (clays, zeolites), Synthetic.
    \item \textbf{Organic:} Natural (peat, coal), Synthetic (resins).
\end{itemize}
\textbf{Zeolites General Formula:} $(\text{Me}^{2+}, \text{Me}_2^+)O \cdot \text{Al}_2\text{O}_3 \cdot n\text{SiO}_2 \cdot m\text{H}_2\text{O}$

\textbf{Ion Exchange Equilibrium:}
\begin{equation}
    \text{RM}_2 + \text{M}_1\text{X} \rightleftharpoons \text{RM}_1 + \text{M}_2\text{X}
\end{equation}
Reversible process, follows the Law of Mass Action.
\vspace{0.2cm}